\SourcePage{144}{147}
\section{Expansion of a plane wave}

Consider a free particle moving with definite momentum $p=k\hbar$ along the
positive $z$ axis. Its wave function is
\[
 \psi=\operatorname{const}\cdot e^{ikz}.
\]
Expand it in the free-motion functions $\psi_{klm}$ with definite angular
momentum. Since $E=k^2\hbar^2/2m$ is definite, only functions with the same
$k$ enter. Since $e^{ikz}$ is axially symmetric about $z$, only functions
independent of $\varphi$, with $m=0$, can occur. Thus
\SourcePage{145}{148}
\[
 e^{ikz}=\sum_{l=0}^{\infty}a_l\psi_{kl0}
 =\sum_{l=0}^{\infty}a_lR_{kl}Y_{l0}.
\]
Substitution of (28.8) and (33.9) gives
\[
 e^{ikz}=\sum_{l=0}^{\infty}C_lP_l(\cos\theta)
 \left(\frac rk\right)^l
 \left(\frac1r\frac d{dr}\right)^l\frac{\sin kr}{kr}
 \qquad(z=r\cos\theta).
\]
Determine $C_l$ by comparing coefficients of $(r\cos\theta)^n$ in the
power-series expansions. On the right this term occurs only in the $n$th
summand: for $l>n$ the radial expansion begins with higher powers of $r$,
while for $n>l$, $P_l(\cos\theta)$ has lower powers of $\cos\theta$. The
coefficient of $\cos^l\theta$ in $P_l(\cos\theta)$ is
$(2l)!/2^l(l!)^2$ (see (c.1)). By (33.13), the relevant term is
\[
 (-1)^lC_l\frac{(2l)!(kr\cos\theta)^l}
 {2^l(l!)^2\,1\cdot3\ldots(2l+1)}.
\]
The corresponding term in $\exp(ikr\cos\theta)$ is
\[
 \frac{(ikr\cos\theta)^l}{l!}.
\]
Hence $C_l=(-i)^l(2l+1)$, and
\begin{equation}
 e^{ikz}=\sum_{l=0}^{\infty}(-i)^l(2l+1)P_l(\cos\theta)
 \left(\frac rk\right)^l
 \left(\frac1r\frac d{dr}\right)^l\frac{\sin kr}{kr}.
 \tag{34.1}
\end{equation}
At large distances,
\begin{equation}
 e^{ikz}\approx\frac1{kr}\sum_{l=0}^{\infty}i^l(2l+1)
 P_l(\cos\theta)\sin\left(kr-\frac{l\pi}{2}\right).
 \tag{34.2}
\end{equation}
In (34.1), $z$ is along the wave vector $\mathbf k$. The expansion also has a
\SourcePage{146}{149}
coordinate-independent form. The addition theorem for spherical functions
(c.11) expresses $P_l(\cos\theta)$ through spherical functions of the
directions of $\mathbf k$ and $\mathbf r$:
\begin{equation}
 e^{i\mathbf k\mathbf r}=4\pi\sum_{l=0}^{\infty}\sum_{m=-l}^{l}
 i^lj_l(kr)Y_{lm}^*\left(\frac{\mathbf k}{k}\right)
 Y_{lm}\left(\frac{\mathbf r}{r}\right).
 \tag{34.3}
\end{equation}
Since $j_l(kr)$ depends only on $kr$, the symmetry in $\mathbf k$ and
$\mathbf r$ is explicit; which spherical function is conjugated is immaterial.

Normalize $e^{ikz}$ to unit probability-current density: one particle crosses
unit area per unit time. The function is
\begin{equation}
 \psi=\frac1{\sqrt v}e^{ikz}=\sqrt{\frac m{k\hbar}}e^{ikz}
 \tag{34.4}
\end{equation}
($v$ is the speed; see (19.7)). Multiplying (34.1) by $\sqrt{m/k\hbar}$ and
using $\psi_{klm}^{\pm}=R_{kl}^{\pm}Y_{lm}$ gives
\[
 \psi=\sum_{l=0}^{\infty}\sqrt{\pi(2l+1)}\,
 \frac1{ik\hbar}(\psi_{kl0}^{+}-\psi_{kl0}^{-}).
\]
The squared coefficient of $\psi_{kl0}^{-}$ (or $\psi_{kl0}^{+}$) gives the
probability that a particle incident on the center (or leaving it) has angular
momentum $l$. Since the incident flux density is unity, this probability has
dimensions of area and can be interpreted as the target area in the $xy$
plane for particles with angular momentum $l$. Denoting it by $\sigma_l$,
\begin{equation}
 \sigma_l=\frac\pi{k^2}(2l+1).
 \tag{34.5}
\end{equation}
\SourcePage{147}{150}
For large $l$, summing over a range $\Delta l$, where
$1\ll\Delta l\ll l$, gives
\[
 \sum_{\Delta l}\sigma_l\approx\frac\pi{k^2}\cdot2l\Delta l
 =2\pi\frac{l\hbar^2}{p^2}\Delta l.
\]
With the classical relation $\hbar l=\rho p$, where $\rho$ is the impact
parameter, this becomes
\[
 2\pi\rho\Delta\rho,
\]
the classical result. This is not accidental: for large $l$ the motion is
semiclassical (\secref{49}).

\ProblemHeading{} Expand a plane wave in wave functions with definite
projection $m$ of angular momentum on the $y$ axis and definite momentum
projection $p_y$ on that axis.

\SolutionHeading{} Introduce cylindrical coordinates $y,\rho,\varphi$ with
axis $y$. The functions have the form
$Q_m(\rho)e^{im\varphi}e^{ip_yy/\hbar}$. Measuring $\varphi$ from $z$,
\[
 e^{ikz}=e^{ik\rho\cos\varphi}
 =\sum_{m=-\infty}^{\infty}Q_m(\rho)e^{im\varphi}
\]
(here $p_y=0$), whence
\[
 Q_m(\rho)=\frac1{2\pi}\int_0^{2\pi}
 \exp[i(k\rho\cos\varphi-m\varphi)]\,d\varphi=i^mJ_m(k\rho),
\]
where $J_m$ is a Bessel function. For $k\rho\gg1$,
\[
 Q_m(\rho)\approx i^m\sqrt{\frac2{\pi k\rho}}
 \sin\left[k\rho-\frac\pi2\left(m-\frac12\right)\right].
\]
